\section{Reliability, Security, and Governance}
\label{sec:safety}

Safety in an agentic system is not primarily a property of its generated text; it is a property of the transition from evidence to external effect. A chatbot can produce harmful content, but an agent can additionally read private data, admit information into persistent memory, invoke tools, communicate with third parties, and modify authoritative state. A failure becomes consequential when untrusted or unsupported evidence is admitted as state, when data is interpreted as instruction, when a proposed action receives authority beyond the user's purpose, or when an effect is accepted without verification. Risk therefore arises across the entire evidence--state--action path rather than at the model boundary alone.

This view turns safety from a late output filter into a set of enforceable representation and control boundaries. Retrieved content, user instructions, tool results, memories, plans, action proposals, authorizations, and execution receipts should remain distinguishable through typed records with provenance and scope. The runtime can then apply least privilege, policy checks, confirmation, and effect verification at the transition where each control matters. OWASP's excessive-agency category identifies excessive functionality, permissions, and autonomy as distinct causes of damaging action~\cite{owasp2025llmtop10}; the architectural response is correspondingly to bound all three rather than ask the language model to compensate for an overpowered runtime. Safety evidence must therefore show both that hazardous transitions are blocked and that legitimate work remains possible.

\subsection{Threat model}

A threat model makes those boundaries explicit. It should identify protected assets, actors, trust zones, admissible information flows, granted authorities, permitted effects, and recovery obligations. Relevant assets include credentials, private documents, user identity, financial resources, source code, memory stores, tool registries, decision records, and the integrity of external systems. Actors may include malicious users, compromised documents or websites, untrusted tool servers, peer agents, insiders, and accidental misconfiguration; the model and planner should also be treated as fallible components rather than trusted policy authorities.

For each possible state transition, the threat model should ask: who introduced the evidence; how its origin and integrity are represented; which principal and task authorize its use; what data the transition may read or disclose; which external effects it may produce; how those effects are verified; and how the system contains or reverses failure. This converts a generic asset inventory into testable security contracts. For example, a retrieved email may be admissible as evidence about a user's request while remaining inadmissible as authorization to send data, change memory, or invoke a financial tool.

The empirical literature covers complementary parts of this contract rather than one universal notion of security. InjecAgent contains 1,054 indirect prompt injection cases spanning 17 user tools and 62 attacker tools; among its tested configurations, a ReAct-prompted GPT-4 agent executed the attack in 24\% of cases~\cite{zhan2024injecagent}. Agent Security Bench broadens coverage to ten scenarios, more than 400 tools and tasks, and 27 attack or defense methods; it reports a highest average attack success rate above 84.3\% across its evaluated configurations~\cite{zhang2024agentsecuritybench}. ST-WebAgentBench instead attaches 3,057 safety and trustworthiness policies to 375 web tasks and measures six dimensions, explicitly coupling constraint satisfaction to useful task completion~\cite{levy2026stwebagentbench}. These results are not directly rankable because their actors, actions, policies, and acceptance rules differ. Together they support a narrower systems conclusion: as connectivity, persistence, and authority expand, every newly reachable transition requires a corresponding evidence label, policy gate, and effect test. Threat knowledge bases such as MITRE ATLAS help enumerate adversarial tactics, but they must be instantiated in such task-level contracts rather than treated as performance evidence~\cite{mitre2026atlas}.

\subsection{Normative guidance and empirical security evidence}

Agent assurance requires an evidence stack. NIST frameworks organize governance and risk work~\cite{nist2023airmf,nist2024genaiprofile}; OWASP and MITRE ATLAS enumerate vulnerability classes and tactics~\cite{owasp2025llmtop10,mitre2026atlas}; controlled benchmarks test concrete systems under authored threats. None substitutes for another: a framework does not demonstrate prompt injection resistance, a taxonomy does not estimate control effectiveness, and a benchmark pass says nothing about untested threats.

The layers become coherent as an assurance case linking a scoped claim to a threat, control, test contract, versioned evidence, verdict, owner, and residual uncertainty. Stable identifiers and lifecycle state make applicability queryable when models, tools, data, or permissions change. \Cref{tab:safety-evidence} therefore compares what each source supports and where its inference must stop.

\begin{table}[H]
\centering
\caption{Evidence stack for agent assurance. The sources are complementary rather than substitutable; the final column states the supported use and its limit.}
\label{tab:safety-evidence}
\small
\surveytableformat
\begin{tabularx}{\textwidth}{L{2.8cm}L{2.5cm}L{4.15cm}Y}
\toprule
\rowcolor{surveyBlue!28}\textbf{Source} & \textbf{Instrument} & \textbf{What it contributes} & \textbf{Supported use and limit} \\
\midrule
NIST AI RMF~\cite{nist2023airmf} & Normative framework & Govern, map, measure, and manage outcomes across the lifecycle & Assign ownership and structure risk work; does not validate a particular runtime or control \\
NIST GenAI Profile~\cite{nist2024genaiprofile} & Cross-sector risk profile & Generative AI risks and suggested lifecycle actions & Extend the control portfolio to model, data, content, and human system risks; not an agent test \\
OWASP LLM Top 10~\cite{owasp2025llmtop10} & Vulnerability taxonomy & Application risks including prompt injection, improper output handling, and excessive agency & Build coverage and design checklists; does not measure prevalence or mitigation effectiveness \\
MITRE ATLAS~\cite{mitre2026atlas} & Living adversary catalog & Tactics, techniques, mitigations, and case study vocabulary for AI-enabled systems & Map threats, red-team cases, detection, and incidents; coverage evolves and is not a system score \\
InjecAgent~\cite{zhan2024injecagent} & Controlled benchmark & Indirect prompt injection through tool-integrated content & Estimate attack and defense behavior for covered tools and attacks; not general agent security \\
Agent Security Bench~\cite{zhang2024agentsecuritybench} & Controlled benchmark & Multiple attack and defense classes across tool-rich settings & Compare resilience under its configurations; claims remain bound to its scenarios, agents, and metrics \\
ST-WebAgent\-Bench\newline\cite{levy2026stwebagentbench} & Policy-and-utility benchmark & Safety constraints and useful completion in stateful web tasks & Test policy compliance without rewarding blanket refusal; scope remains the authored web workflows \\
\bottomrule
\end{tabularx}
\end{table}

The practical implication is a claim-to-evidence chain. NIST can assign responsibility, OWASP and ATLAS map threat coverage, and benchmarks exercise specified transitions. The result supports only the observed controls and effects; residual threats, untested permissions, and evaluator blind spots remain open claims.

\subsection{Prompt injection and instruction hierarchy}

Prompt injection is representation collapse: source, type, and trust are flattened into text, allowing data to be interpreted as control. Direct attacks supply competing instructions; indirect attacks place them in later retrieved content~\cite{greshake2023promptinjection}. The consequential failure is not textual disobedience but lower-trust evidence altering an objective, invoking authority, or redirecting an effect.

InjecAgent's 1,054 cases and AgentDojo's 629 security tests expose complementary tool-mediated attacks~\cite{zhan2024injecagent,debenedetti2024agentdojo}. Their outcomes depend on prompts, tools, permissions, task state, and evaluator policy rather than the model alone. Adaptive attacks exceeded 50\% success against all eight studied defenses~\cite{zhan2025adaptiveattacks}, showing that a fixed prompt suite is a bounded test contract, not a durable security boundary.

Instruction hierarchy and StruQ restore separation between higher-trust instructions and lower-trust data and report gains with limited utility loss~\cite{wallace2024instructionhierarchy,chen2025struq}. Their scope is model interpretation. A trust label matters only when created outside attacker control, and ranked instructions do not authenticate a principal or authorize an effect.

A reliable architecture must preserve the trust distinction across the complete evidence-to-action path rather than ask the model to recover it from wording. The resulting controls form a linked chain:
\begin{itemize}
    \item \textbf{Ingest as evidence.} Store external content in typed records carrying origin, principal, channel, retrieval scope, integrity state, and permitted use; sanitization may remove active content but must not be treated as proof of trust.
    \item \textbf{Assemble context by policy.} Keep instructions and evidence in separate fields, compute precedence from authenticated channels rather than self-declared text, and minimize secrets or irrelevant instructions exposed to the model.
    \item \textbf{Constrain plan formation.} Propagate provenance from evidence into claims and proposed arguments; content may support a task but may not silently expand its objective, permissions, recipients, or action space.
    \item \textbf{Authorize outside the model.} Validate every tool request against the authenticated principal, declared purpose, resource scope, current state, and action class. The model proposes an invocation; a deterministic monitor decides whether it is admissible.
    \item \textbf{Separate proposal from effect.} Require explicit confirmation for disclosure, communication, financial commitment, or irreversible change, then bind the approval to the exact arguments and retain an authoritative execution receipt.
    \item \textbf{Evaluate the chain.} Measure benign utility, targeted and adaptive attack success, policy-violating effects, false refusals, and recovery across adversarial documents, compromised tool outputs, and multi-step exfiltration attempts.
\end{itemize}
The controls are complementary. Typed context without external authorization can still turn model error into action; authorization without provenance can approve poisoned arguments; and confirmation without an exact proposal is ambiguous. Prompt injection resistance is therefore the preservation of authority despite malicious strings, with the model treated as interpreter and proposer rather than final security boundary.

\subsection{Tool misuse and confused-deputy failures}

Prompt injection explains how low-trust evidence can influence a decision; the confused-deputy problem explains why that influence becomes a security violation. A deputy is confused when it exercises authority that it legitimately holds, but for a requester or purpose that was never entitled to use it~\cite{hardy1988confuseddeputy}. In an agent, the requester may be a user, a retrieved document, a remote agent, or a tool response. If a document can cause an executor holding the user's cloud credential to upload a private file, the failure is not simply incorrect tool selection. Authority from one principal has been redirected by evidence from another.

Preventing redirection requires separating tool description, principal delegation, concrete action proposal, authorization decision, and execution receipt. Delegation defines purpose and scope; the proposal preserves argument provenance; the decision binds one version and current state; the receipt records the effect. A JSON schema validates shape but cannot replace these authority-bearing records.

Benchmarks reveal complementary parts of the problem. ToolEmu uses an LM-emulated sandbox to exercise 144 cases across nine risk types and 36 toolkits, making hazardous trajectories observable without executing them against real systems~\cite{ruan2024toolemu}. Agent Security Bench broadens coverage to ten scenarios, ten agents, more than 400 tools, 27 attack or defense methods, and seven evaluation metrics~\cite{zhang2024agentsecuritybench}. These suites identify risky behaviors and comparative weaknesses, but neither a simulated consequence nor a benchmark verdict grants deployment authority. Conversely, step-level monitoring can improve the transition immediately before execution: ToolSafe reports a 65\% average reduction in harmful tool invocations and roughly a 10\% improvement in benign completion under prompt injection~\cite{mou2026toolsafe}. Such learned guards are useful evidence and defense layers, yet their predictions remain inputs to authorization rather than authorization itself.

Least privilege means narrow executable authority, not fewer tool names. Separate read from write, scope credentials to task and resource, mediate calls after arguments are known, and separate high-impact proposal from commit. Values derived from untrusted content need provenance-aware flow controls, and policy must be rechecked against current state at dispatch.

CaMeL illustrates the architectural alternative to relying on model restraint. It extracts control and data flows from a trusted query, prevents untrusted data from changing program flow, and uses capabilities to block unauthorized exfiltration; within its stated threat model, it completes 67\% of AgentDojo tasks with provable security~\cite{debenedetti2025camel}. The result is important precisely because it separates two questions that ordinary agent loops conflate: whether the model can interpret evidence and whether the resulting information flow is authorized. Its remaining utility gap also makes the trade-off visible rather than hiding it behind a single success score.

Tool output remains untrusted evidence after success and must be schema-, size-, provenance-, and content-checked. The governing invariant is \emph{authority non-amplification}: evidence may refine an authorized action but cannot mint authority, change its principal, or broaden its purpose. The final effect must remain a subset of explicit delegation.

\subsection{Memory poisoning and privacy}

Persistent memory amplifies trust across later tasks, tools, and users. The failure occurs when storage silently promotes an untrusted message to fact, inference to preference, or trajectory to procedure. MemoryBank makes long-term updating and forgetting explicit~\cite{zhong2024memorybank}, but persistence still requires distinct types and writers for observations, facts, preferences, summaries, procedures, policies, and authorization artifacts.

Memory records therefore separate payload from immutable identity, type, subject, provenance, validity, sensitivity, purpose, scope, retention, and supersession metadata. Quoted text cannot declare its own authority; summaries retain lineage, corrections supersede, and revocation propagates without preserving unnecessary sensitive payload.

Capability benchmarks and attack studies expose different inference boundaries. MemBench evaluates factual and reflective memory across participation and observation settings, emphasizing effectiveness, efficiency, and capacity~\cite{tan2025membench}; those measurements do not test whether a remembered item is authorized or adversarial. AgentPoison demonstrates the integrity gap directly: poisoning less than 0.1\% of a long-term memory or RAG knowledge base produced average attack success rates of at least 80\% across three agent settings while changing benign performance by at most 1\%~\cite{chen2024agentpoison}. MINJA removes the assumption of direct storage access and injects malicious records through ordinary query interactions and observed outputs~\cite{dong2025minja}. Together, these results show that the write path and the retrieval policy form a joint attack surface; protecting only the database interface leaves model-mediated memory creation exposed.

A secure lifecycle uses separate write and read gates. Admission authenticates the writer, classifies type, checks conflict and sensitivity, and assigns scope and expiry. Retrieval applies access, purpose, validity, and revocation filters before similarity. Returned memory enters planning as typed evidence, not privileged instruction; current task policy decides its influence.

Derived representations inherit privacy obligations. Embedding inversion recovered 92\% of tested 32-token inputs and reconstructed personal information from clinical notes~\cite{morris2023embeddinginversion}, so vectors, summaries, entities, caches, and traces are not anonymized substitutes. Access, minimization, retention, and deletion must propagate across every derivative.

Evaluation should test adversarial writes, stale and conflicting records, poisoned summaries, scope leakage, inversion, trace disclosure, and deletion. Utility, integrity, privacy, update correctness, and action safety remain separate outcomes. The invariant is \emph{scoped persistence}: memory cannot outlive its validity, escape its audience or purpose, or grant new authority.

\subsection{Protocol and supply chain security}

Interoperability is not only a connectivity problem; it is an authority composition problem. MCP standardizes the discovery and invocation of resources, prompts, and tools, while A2A uses Agent Cards to advertise identities, endpoints, skills, and authentication requirements~\cite{mcp2025spec,a2a2026spec}. These records make heterogeneous components legible to one another, but they remain claims supplied by another administrative domain. Indeed, the MCP specification explicitly treats tool descriptions and annotations as untrusted unless they come from a trusted server~\cite{mcp2025spec}. A syntactically valid schema or signed Agent Card can establish what an endpoint says about itself; neither establishes that the advertised behavior is safe for the present user, task, or data. If discovery metadata is allowed to confer authority directly, adding a connector silently enlarges the system's privileges and transitive trust.

The control plane therefore maintains a connector manifest binding reviewed identity, version, capabilities, scopes, data flows, side effects, reach, approval, validity, revocation, and tests. Its lifecycle is \emph{discover--review--pin--activate--invoke--monitor--revoke}. Material changes to endpoint, schema, scope, or side effect trigger disablement or reapproval and preserve an incident inventory.

Authentication protects the boundary only with delegation. MCP requires audience-bound tokens and rejects token pass-through~\cite{mcp2025spec}; A2A advertises authentication separately from credentials~\cite{a2a2026spec}. Remote agents should receive minimal task context and narrow, audience-bound authority, never ambient identity or memory. Returned artifacts remain untrusted evidence.

This separation is necessary because protocol-level attacks target the gap between description and behavior. MCP Security Bench operationalizes attacks including name collisions, malicious instructions in tool descriptions, out-of-scope parameter requests, user-impersonating responses, retrieval injection, and transfers between malicious tools; across more than 400 real tools, the study finds that stronger tool use does not by itself imply stronger resistance~\cite{zhang2025mcpsecuritybench}. Static allow lists are therefore insufficient: evaluation should exercise malicious servers, tool shadowing, metadata and schema drift after approval, credential audience or scope substitution, compromised updates, exfiltration through tool results, cross-agent privilege escalation, and revocation during an active task. The outcome of interest is not merely whether the protocol exchange succeeds, but whether the composed workflow remains inside the authority recorded in the manifest.

in-toto and SLSA can link a deployed connector to reviewed production provenance~\cite{torresarias2019intoto,slsa2026provenance}, but do not prove benign behavior or task authority. The invariant is \emph{bounded composition}: protocol metadata describes capabilities; only local policy against pinned identity, version, and delegation may grant them.

\subsection{Side effects and transactional safety}

Authorization does not make an action operationally correct. Once a request crosses a tool boundary, the response channel and the effect channel can diverge: a timeout may arrive after a payment was accepted, a successful response may precede eventual visibility, a concurrent actor may invalidate a checked balance or file version, and a multi-system workflow may stop after only some writes. Blind retry can then duplicate an effect, while trusting the response can produce a false claim of completion. Transactional safety therefore asks a second question after ``may this action run?'': what exact state transition occurred, under which preconditions, and what recovery obligation remains?

The first control is to classify actions by impact and reversibility because different effects require different commit protocols. A useful taxonomy is:
\begin{itemize}
    \item \textbf{read-only}: inspect data without modifying external state;
    \item \textbf{reversible write}: create a draft, branch, or staged change that can be undone;
    \item \textbf{externally visible communication}: send a message, publish content, or contact a third party;
    \item \textbf{financial or legal commitment}: transfer value, accept terms, or submit an official form;
    \item \textbf{destructive or safety-critical action}: delete data, change access, operate equipment, or alter production systems.
\end{itemize}
This classification is not merely a label for an approval screen. Read-only calls can often be repeated, although their freshness and hidden service-side effects still matter. Reversible writes can execute into a draft, branch, escrow, or staging area before promotion. Communication, financial commitment, and destructive change instead require a narrow authorization bound to the final target and arguments, stronger preconditions, and an explicit commit point; a later apology, refund, or restoration is not equivalent to preventing the original effect. The required authorization and verification should therefore increase with impact, irreversibility, uncertainty, and the number of external systems involved.

Each intended effect should be an \emph{action transaction} linking the exact proposal, principal, target, arguments, preconditions, effect class, idempotency, approval, attempts, postconditions, compensation, and receipts. Its lifecycle separates proposal, authorization, dispatch, commit, verification, abortion, compensation, and reconciliation. Only the runtime advances the record from authenticated policy and external evidence; the goal is durable orchestration, not fictional atomicity across unrelated APIs.

Tool adapters make this record executable by declaring preconditions, postconditions, effect class, retry semantics, and available compensation. ToolGate illustrates the local contract boundary: it maintains typed symbolic state and uses Hoare-style preconditions to gate invocation and postconditions to decide whether a result may update trusted state~\cite{liu2026toolgate}. GoEX emphasizes a complementary boundary, using undo and damage confinement to support post-facto validation of actions that have already executed~\cite{patil2024goex}. The former prevents an invalid transition from being committed to the agent's state; the latter limits or reverses damage after execution. Neither alone establishes atomicity across independent services, so contract validation, containment, and workflow recovery should be treated as complementary controls rather than competing definitions of reliability.

Before dispatch, the executor revalidates proposal, authority, arguments, and resource version. Retryable mutations use a stable idempotency key enforced by the authoritative service; HTTP warns against automatic retry of non-idempotent requests without additional knowledge~\cite{fielding2022httpsemantics}. After ambiguity, reconcile state before retrying. In controlled simulations, verify-before-retry raised success from about 58\% to 80\% and reduced duplicates from about 42\% to 20\%~\cite{mansoor2026verifiedtoolcalls}, illustrating why recovery mechanisms require separate measurement.

When no common atomic commit exists, sagas decompose work into local transactions and explicit compensation~\cite{garciamolina1987sagas}. Compensation is semantic repair, not time reversal. Atomix delays bufferable effects and tracks compensation; under its stated 0.30 fault setting, the full configuration achieved 57\% clean success versus 0--7\% for several baselines, while checkpoint replay was statistically tied at 53\%~\cite{mohammadi2026atomix}. The comparison supports structured recovery without making one runtime universal.

Evaluation should inject failures at every effect boundary: after dispatch but before response, during delayed visibility, on duplicate delivery, under stale preconditions and concurrent writes, after approval revocation, between steps of a multi-system workflow, and while compensation itself fails. Useful outcomes include unauthorized commit rate, duplicate effect rate, false completion rate, residual state divergence, recovery latency, uncompensated harm, and benign completion; collapsing them into task success hides the reason a workflow appears reliable. The governing invariant is \emph{verified commit}: an action may be reported as complete only when the exact authorized proposal was executed against valid preconditions and an authoritative postcondition or receipt confirms the resulting state. Ambiguity remains a reconciliation obligation, not a success state.

\subsection{Human oversight}

Human oversight is not the presence of a person somewhere in the workflow; it is a transfer of decision authority at a specific state transition. A reviewer must have the appropriate role, receive the exact proposal and relevant evidence, understand why intervention is requested, have enough time to act, and possess an enforceable means to approve, modify, reject, pause, or take over. A notification after an irreversible effect, or an approval button whose decision is not bound to the executed version, provides visibility without control. The allocation of functions between people and automation is therefore an architectural choice rather than a generic requirement to keep a human ``in the loop''~\cite{parasuraman2000levels}.

Oversight should be placed where human judgment changes the risk of the next transition. Useful modes include co-planning before a plan becomes executable, checkpoint review before a long-running workflow crosses phases, action approval immediately before a high-impact commit, co-tasking or takeover when the system lacks information or control, and post-effect verification when an effect is reversible and damage is bounded. Magentic-UI instantiates these interaction modes through editable plans, co-tasking, action guards, answer verification, multitasking, and saved workflows~\cite{mozannar2025magenticui}. Its autonomous, simulated user, qualitative, and targeted safety evaluations make the design space concrete, but they do not establish that every mechanism improves every task. Placement should instead depend on consequence, reversibility, evidence quality, verification strength, time sensitivity, and reviewer capacity. Low-risk, objectively verified transitions may use exception-based review; external communication, financial commitment, permission change, or destructive action should retain a gate at the final commitment boundary.

A typed \emph{review request} binds one proposal and policy trigger to affected resources, evidence, uncertainty, alternatives, constraints, expected effects, reviewer role, deadline, safe default, and recovery. Interfaces may summarize, but should expose material conflict and the exact change under review. This turns ``ask the user'' into a routable control object rather than another chat turn.

More review is not automatically safer. Imperfect aids produce omission, commission, and complacency errors~\cite{parasuraman2010complacency}, while passive automation can reduce situation awareness and slow takeover~\cite{endsley1995outoftheloop}. Interfaces should use selective gates, clear consequence and provenance displays, easy pause and takeover, and workload limits rather than approval fatigue. Risk and reversibility matter more than model confidence alone.

Escalation is a routing policy. Learning-to-defer models machine and expert jointly~\cite{mozannar2020defer}; practical triggers include missing authority, conflict, novelty, weak verification, non-progress, or effects beyond delegation. The system must route to relevant authority and distinguish requests for information, judgment, approval, or execution. Field evidence finds different outcomes for technical and emotional escalations and an effect of timing~\cite{wang2026agentichitl}, so type and timing belong in evaluation.

The decision record binds reviewer, role, rationale, conditions, expiry, and one proposal version. Modification creates a new proposal that must be rechecked; approval cannot float to a changed action. The receipt shows what occurred. Human correction remains episode-scoped unless separately reviewed for promotion into memory, procedure, or policy.

Evaluation must cover the combined human--agent system rather than report model accuracy plus the existence of a gate. Comparisons should include no-review, blanket review, and risk-triggered-review conditions and separately measure hazardous approvals, correct interventions, unnecessary escalations, reviewer workload, decision latency, time to takeover, recovery quality, residual harm, and benign completion. Adversarial cases should present stale proposals, misleading summaries, omitted conflicts, time pressure, sequences of harmless requests before a dangerous one, and edits made after approval. The governing invariant is \emph{meaningful intervention}: a human decision counts as a control only when an authorized reviewer receives the correct version and material evidence with enough time and agency to alter the outcome, and the resulting effect remains traceable to that decision.

\subsection{Observability, audit, and incident response}

Observability is not a larger transcript or a dashboard of token counts. It is the preservation of enough typed relationships to reconstruct why a state transition occurred, which authority it used, what evidence it consumed, and what effect it produced. This distinction matters because an agentic run crosses model calls, retrieval services, memory, policy engines, tools, human gates, and external systems. Each component may emit a locally correct log while the end-to-end account remains ambiguous. In delegated or multi-agent execution, the same low-level actions can even be compatible with different assignments of authority unless delegation context is bound to the events when they occur~\cite{mishra2026delegatedobservability}. Observability therefore continues the paper's central argument: evidence becomes governable action only when the system preserves the representations that connect proposal, authorization, execution, and result.

The basic representation is an append-only \emph{trace event} with causal identity, actor and delegation, component version, artifact links, decision and receipt references, status, error, cost, integrity, and retention. Parent and correlation edges connect decomposition, evidence use, handoffs, retries, compensation, and asynchronous work. The result is a causal graph rather than interleaved text.

At minimum, the linked spans and records should cover:
\begin{itemize}
    \item the task, principal, purpose, constraints, authorization scope, and delegation chain;
    \item model, prompt, policy, planner, evaluator, memory, and tool versions;
    \item retrieved evidence and memory items with identifiers, versions, provenance, and applicability;
    \item plans, state transitions, routing decisions, and the alternatives rejected at consequential branches;
    \item tool proposals, validated arguments, precondition checks, requests, responses, retries, and errors;
    \item guardrail and policy decisions, including the rule, result, and policy version;
    \item human review requests, approvals, modifications, overrides, expiries, and the proposal version reviewed;
    \item external side effects, compensation attempts, execution receipts, and final authoritative state; and
    \item token, latency, compute, monetary, reviewer, and recovery costs.
\end{itemize}
The AgentOps taxonomy similarly treats observability as lifecycle coverage across agent artifacts and their associated data rather than as model call logging alone~\cite{dong2024agentops}. OpenTelemetry's emerging GenAI conventions provide a useful interoperability layer for model, agent, and tool spans~\cite{opentelemetry2026genai}; however, a transport schema does not by itself capture the survey's semantic objects: evidence records, memory provenance, proposal versions, authority transfers, and verified effects. Those relationships should remain explicit even when exported through a common telemetry substrate.

Instrumentation must also separate what the agent \emph{represented} from what the environment \emph{did}. Application-level spans can expose goals, plan revisions, retrieval, routing, policy checks, and structured decision summaries. Boundary and system-level telemetry can independently observe network requests, file and process activity, resource use, and external commits. AgentSight illustrates the value of correlating these layers: its boundary tracing design joins semantic intent to system effects, reports less than 3\% performance overhead, and detects prompt injection behavior, wasteful reasoning loops, and multi-agent coordination bottlenecks in its evaluated settings~\cite{zheng2025agentsight}. The comparison is architectural rather than universal: application instrumentation is semantically rich but framework-dependent, whereas boundary tracing is harder for the agent to bypass but cannot recover evidence or authority relationships that were never encoded. High-risk systems need both views and should treat the external receipt as authoritative when the two disagree.

TraceElephant reports up to 76.5\% better failure attribution from full than partial traces~\cite{chen2026traceelephant}, supporting context around decisive steps rather than permanent capture of every prompt. Structural metadata and sensitive content need separate access and retention; secrets should become scoped references. Nominal runs may be sampled, but policy exceptions, overrides, failures, and effects require complete structural records.

Replay is useful only when its semantics are explicit. A \emph{trace replay} renders the recorded execution for inspection; a \emph{decision replay} reruns selected model or policy steps against frozen inputs; and an \emph{effect replay} invokes tools and is therefore unsafe unless performed with mocks, recorded responses, or an isolated test environment. The trace should record which observations were captured, which dependencies were substituted, and where a replay diverged. Because remote data, model behavior, credentials, and time-dependent preconditions change, exact reenactment is often impossible. The practical goal is narrower: localize the first material divergence, test whether a proposed policy or component change would alter it, and preserve the result as a regression case without repeating the original side effect.

Incident response should consume traces through a state machine aligned with preparation, detection, response, and recovery~\cite{nist2025incidentresponse}. A typed incident links detection and severity to affected principals, components, credentials, resources, evidence, containment owner, remediation, and recovery checks. Operators need scoped controls to pause, revoke, quarantine, disable, enumerate, and reconcile; global shutdown is a last resort.

Detection, attribution, and recovery are distinct. AgentDebugX operationalizes Detect--Attribute--Recover--Rerun; in its reported setting, DeepDebug reached 28.8\% exact attribution versus 21.7\% for a single-pass baseline and repaired 13 of 73 GAIA failures versus 4--6 for three alternatives~\cite{zhu2026agentdebugx}. The result is configuration-specific but shows that exposing a trace is not the same as diagnosing or repairing failure.

Evaluation should compare final output inspection, flat logs, typed causal traces, and typed traces augmented with boundary telemetry. Useful measures include trace coverage at consequential transitions, causal link accuracy, failure attribution accuracy and time, false attribution, affected episode recall, containment latency, unauthorized post-containment effects, evidence preservation failures, recovery correctness, residual state divergence, and recurrence on prior incidents. The governing invariant is \emph{forensic closure}: an incident is not closed until the responsible authority can identify the affected executions and resources, reconstruct the material evidence--decision--action path, contain further authority, reconcile authoritative state, and bind corrective changes to regression evidence. Observability then becomes the evidence infrastructure for accountability, not a decorative view of agent activity.

\subsection{Governance across the lifecycle}

Governance connects technical controls to an accountable authority that can permit, constrain, change, suspend, and retire the system. The object being governed is therefore not an isolated model. It is a versioned deployment configuration comprising the model and prompts, planners and evaluators, retrieval and memory policies, tools and credentials, human review routes, runtime policies, user population, operating environment, and allowed effect classes. The NIST AI RMF's govern--map--measure--manage structure provides a useful lifecycle scaffold, while its generative AI profile adds model-specific risk considerations~\cite{nist2023airmf,nist2024genaiprofile}. For agentic systems, however, those functions must be bound to the concrete configuration that can act. A research assistant with read-only access and an operational agent that can modify customer accounts may use the same model while requiring different owners, evidence, controls, and release authority.

The persistent representation for this work is a \emph{system governance record}. It should identify the system and configuration version; accountable owner, risk owner, release authority, and operators; intended and prohibited uses; principals, user populations, and affected parties; autonomy and action boundaries; models, prompts, data flows, memory stores, tools, credentials, policies, and dependencies; a risk register with hazards, controls, residual risk, and acceptance criteria; evaluation and incident evidence; monitoring obligations and thresholds; open exceptions and their expiry; release state and decision; and rollback and decommissioning plans. Links to the evidence, review, trace, incident, and receipt objects defined earlier keep the record compact without severing provenance. A governance claim such as ``financial actions remain bounded'' should point to the policy configuration, adversarial tests, approval design, and operational monitors that support it, not merely to a policy document.

This claim--evidence structure is an \emph{assurance case}: a reasoned argument that identified properties hold for a defined system in a defined context. Recent work on AI safety cases distinguishes assertion-, constraint-, and capability-based claims; demonstrative, comparative, causal, risk-based, and normative arguments; and multiple evidence families, while emphasizing that prompt, scaffold, and deployment changes can invalidate an earlier case~\cite{lee2026safetycases}. The useful unit is therefore a versioned graph whose claims declare scope, assumptions, evidence, owner, review status, and expiry. Contradictory or missing evidence remains visible, and a change impact edge identifies which claims must be reopened when a component, permission, population, or environment changes.

Documentation is necessary but does not govern by itself. A review of 37 documentation frameworks and 22 empirical studies found plausible governance benefits, including risk communication, collaboration, and ethical deliberation, but also weak incentives, limited resources, organizational barriers, and poor integration with existing workflows~\cite{winecoff2024documentation}. A separate analysis of 32,111 Hugging Face model documents found uneven informativeness, with limitations and evaluation among the least completed sections~\cite{liang2024modelcards}. These studies concern broader AI documentation rather than agent releases, but their comparison exposes a general failure mode: a static card can be complete enough to publish and still be stale, disconnected from an executable configuration, or unable to stop a release. The governance record should therefore be generated from and checked against version-controlled manifests, test results, policy state, and deployment telemetry, with human judgment preserved for risk acceptance and contested trade-offs.

Release should be a state machine rather than a one-time approval. A practical progression is \emph{design--sandbox--shadow--read-only--limited action--bounded production}, with alternative \emph{paused}, \emph{rolled back}, and \emph{decommissioned} states. Each promotion record should bind the exact configuration hash to an owner, deployment population, permitted actions and resources, evaluation results, unresolved limitations, monitoring thresholds, time window, rollback trigger, and signed decision. Evidence requirements should rise with impact, irreversibility, scale, uncertainty, and exposure to adversarial inputs. Runtime enforcement must then make the recorded boundary real. In Aegis's reported sandbox evaluation, prompt policy conditioning produced 79 risky comparator path leakage rows across 6,300 rows, whereas 2,100 governed rows produced zero governed mock tool applications and zero governed risky side effect completions; all attempted governed actions retained server-resolved provenance~\cite{mazzocchetti2026aegis}. The authors explicitly limit the claim to that corpus, but the comparison illustrates why a release policy that exists only in prompts is weaker than a policy bound to the dispatch boundary.

Post-deployment monitoring tests whether the evidence supporting release remains valid. It should derive signals from named hazards and claims rather than accumulate generic dashboards: task and verification failures test capability claims; unauthorized or blocked action attempts test authority boundaries; escalations and reviewer workload test the human control design; retrieval conflicts and stale memory use test evidence quality; near misses, incidents, and complaints test the risk model; and drift in users, tools, data, or environment tests the deployment scope. NIST AI 800-4 organizes monitoring challenges across functionality, operations, human factors, security, content, and impacts, drawing on practitioner workshops and a literature review; it also reports fragmented logging, immature methods and information sharing, and difficulty scaling human monitoring~\cite{rao2026monitoring}. These findings support broad coverage while cautioning that current monitoring practice is not yet a validated recipe. Every alert should therefore identify the governance claim it threatens, the affected configuration and population, the responsible owner, and the transition it may trigger.

Organizational design must match this continuous technical loop. Product and system owners define purpose and operation; tool, data, and security owners control dependencies and access; a risk owner maintains the hazard and residual risk view; an independent assessor challenges material claims; a release authority accepts or rejects the remaining risk; and operations and incident roles can pause or roll back the deployment. One person may hold several roles in a small organization, but the decision record should make conflicts and authority visible. Mixed-methods research on public sector agent oversight, combining a literature review with interviews of civil servants in AI-related roles, identifies continuous oversight, integration of governance with operations, and interdepartmental coordination as intensified challenges~\cite{schmitz2025agentoversight}. The evidence is sector-specific, yet it reinforces the architectural requirement that governance decisions, operational controls, and incident evidence share identifiers and escalation paths rather than live in isolated teams.

Change management closes the lifecycle loop. A change to a model, prompt, retrieval corpus, memory policy, evaluator, tool schema, credential, human review rule, or monitoring threshold should produce a configuration diff and reopen the claims that depend on it. The resulting release packet combines targeted tests, normal and adversarial regressions, prior incidents, operational canaries, and an updated residual risk decision. Progressive exposure limits the consequence of uncertainty, but promotion should depend on evidence thresholds rather than elapsed time or absence of complaints. Decommissioning is also a governed transition: stop schedules and queues, revoke credentials and delegations, disable routes, resolve or transfer pending actions, archive or delete memory and traces under retention policy, notify affected users where required, and verify that no orphan authority remains.

Evaluation should compare paper-only approval, static documentation, versioned governance records, and records coupled to release and runtime enforcement. Useful measures include undocumented configuration drift, stale claim rate, control-to-hazard coverage, unauthorized promotion, policy/configuration mismatch, time to identify the responsible owner, rollback latency and correctness, affected population recall, orphan credentials after retirement, recurrence of prior incidents, residual risk exceptions past expiry, and benign completion under the approved boundary. The governing invariant is \emph{authorized evolution}: autonomy may expand only through a recorded promotion of an identified configuration supported by current evidence and accepted residual risk, and it must contract when that evidence, context, or authority no longer holds. Governance then becomes the state machine that keeps evidence and action aligned throughout the system's life.

\subsection{Defense in depth}

Defense in depth is not the accumulation of filters. It is the deliberate allocation of prevention, detection, containment, and recovery responsibilities across distinct trust boundaries, so that one plausible failure cannot create, authorize, execute, and conceal a hazardous effect. This interpretation extends the classic principles of fail-safe defaults, complete mediation, separation of privilege, and least privilege from computer security~\cite{saltzer1975protection} to the evidence--decision--action path of an agent. It also follows system security and system safety engineering in treating assurance as a property of the whole sociotechnical system and its interactions, not a sum of individually reliable components~\cite{ross2022trustworthysystems,leveson2012saferworld}. A prompt instruction, a classifier, a human approval, and a log entry are therefore not interchangeable controls: they occupy different points in the causal path and fail in different ways.

The representation that makes this architecture inspectable is a \emph{defense map}. For each threat or hazard, it records the protected asset or state transition; preventive, detective, containment, and recovery controls; the layer and authority that enforce each control; inputs and trust dependencies; evidence and test identifiers; owner, version, health, and expiry; residual risk and failure assumptions; and escalation and rollback routes. Dependency edges should expose common-mode risk. Two nominally separate checks are not independent if they rely on the same model judgment, prompt-visible labels, tool response, credentials, telemetry channel, or overloaded operator. The defense map thus connects the threat register to the typed evidence, authority, receipt, trace, incident, and governance objects developed throughout this survey.

A defensible agent stack applies that map along the full path:
\begin{enumerate}
    \item at the \emph{evidence boundary}, typed source, trust, provenance, and freshness fields keep untrusted data distinct from instructions and policy;
    \item at the \emph{state boundary}, scoped admission, retrieval, update, retention, and deletion rules prevent memory from silently converting observations into durable authority;
    \item at the \emph{planning boundary}, constrained plan and argument objects bind proposed steps to evidence, assumptions, permissions, budgets, and verification conditions;
    \item at the \emph{authority boundary}, bounded delegation and least privilege credentials limit principals, resources, actions, duration, and downstream redelegation, with complete mediation on every consequential use;
    \item at the \emph{dispatch boundary}, typed schemas, deterministic policy, preconditions, and fail-closed validation decide whether a proposal may become a tool call;
    \item at the \emph{effect boundary}, propose--commit separation, version-bound approval, idempotency keys, receipts, and compensation constrain irreversible or externally visible changes;
    \item at the \emph{isolation boundary}, sandbox, network, filesystem, secret, time, and resource limits contain compromised code, tools, and dependencies;
    \item at the \emph{human boundary}, risk-triggered review, faithful review packets, workload limits, and tested takeover routes reserve judgment for decisions whose residual risk requires it;
    \item at the \emph{detection and recovery boundary}, independent telemetry, objective postconditions, non-progress checks, trace replay, incident response, and rollback expose and contain failures that prevention misses; and
    \item at the \emph{lifecycle boundary}, adaptive attacks, fault injection, regression tests, staged release, monitoring, and decommissioning determine whether the preceding controls remain effective as the system changes.
\end{enumerate}

The ordering matters because later layers should not merely repeat an earlier model judgment. The proposer should not be its sole evaluator; a self-declared ``trusted'' label should not establish provenance; the tool response that claims an effect occurred should not be the only postcondition; and a human gate is not independent if the agent can omit material evidence from the review packet or saturate the reviewer. Logging alone is detection evidence, not prevention or containment. Conversely, a strict preventive gate without receipts, monitoring, or recovery can fail silently when a dependency changes. Defense in depth therefore requires both \emph{coverage} across the causal path and \emph{independence} among the controls protecting the same high-impact transition.

Comparative evidence supports this architectural interpretation while delimiting what any one layer proves. In AgentDojo, an adaptive attacker exceeded 50\% attack success against every one of eight prompt injection defenses studied, showing that fixed prompt-level defenses should not be treated as a complete boundary~\cite{zhan2025adaptiveattacks}. CaMeL instead separates control flow from untrusted data and reports 67\% task completion with provable security under its stated threat model, but its utility gap and assumptions leave roles for other layers~\cite{debenedetti2025camel}. ToolSafe reports an average 65\% reduction in harmful tool invocations and roughly 10\% improvement in benign task completion under prompt injection, yet its learned judgment remains a prediction rather than an execution authority~\cite{mou2026toolsafe}. Aegis's sandbox comparison illustrates the complementary role of an external dispatch gate: prompt policy conditioning leaked 79 risky comparator paths across 6,300 rows, whereas 2,100 governed rows produced no governed mock tool applications or risky effects in that corpus~\cite{mazzocchetti2026aegis}. These results are not directly interchangeable benchmarks, but together they show why instruction robustness, architectural information flow control, learned risk assessment, and deterministic action enforcement should be represented as different defenses rather than collapsed into a single safety score.

Prevention must likewise be connected to diagnosis and recovery. Under a 0.30 injected fault rate, Atomix reports 57\% clean task success versus 0--7\% for several baselines, while checkpoint replay reached a statistically tied 53\%, indicating that structured recovery can preserve progress after faults that a front-end filter cannot prevent~\cite{mohammadi2026atomix}. TraceElephant reports up to 76.5\% better failure attribution from full traces than from partial traces, supporting complete causal evidence for incident analysis without implying that traces themselves block harm~\cite{chen2026traceelephant}. Human review adds a distinct form of judgment, but automation complacency and out-of-the-loop performance problems mean that it should be evaluated as an operational control with measurable information quality, workload, and takeover performance rather than assumed to be a universal backstop~\cite{parasuraman2010complacency,endsley1995outoftheloop}.

Evaluation should compare a single control, staged combinations, and the full mapped stack under the same consequential tasks. Tests should include adaptive prompt injection, poisoned or stale state, compromised tools and monitors, invalid or revoked approval, misleading review packets, telemetry loss, partial execution, recovery failure, and changes that invalidate earlier evidence. Useful measures include hazardous effect escape rate, common-mode bypass rate, hazard-to-control coverage, control independence violations, false refusal, benign completion, reviewer load, containment latency, residual harm, attribution accuracy, recovery correctness, and recurrence after a corrective release. Layer-ablation tests reveal whether an apparent defense is truly independent or simply redundant under ordinary inputs.

The governing invariant is \emph{bounded failure}: for every high-impact transition, no single plausible model, data, tool, policy, operator, or telemetry failure may both produce the hazardous proposal and remove all independent barriers, detection, and recovery; any escaped effect must remain attributable, containable, and recoverable within the declared boundary. Proportionality then becomes precise rather than rhetorical. A lower-impact, reversible action may justify fewer layers, whereas irreversible, high-scale, or rights-affecting actions require stronger coverage and independence. Autonomy should expand only when measured benefit and current evidence justify the residual risk of the complete stack, and should contract when a control, dependency, or operating assumption no longer holds.
